\bookchapter{Глава 38. Пять принципов мышления в ИИ-программировании}

\chapterintro{Если из всей книги вы должны запомнить лишь одну главу, пусть это будет эта. Пять принципов мышления — сердце книги; все последующие главы — контекст, требования, диалог, качество — лишь разворачивают эти пять принципов в разных ситуациях. Не бойтесь: каждый из них прост для понимания, но каждый стоит «вживить в мышцы» с помощью нескольких проектов.}

История начинается занятно: выкладывая свой проект в открытый доступ, Andrej Karpathy приложил файл CLAUDE.md всего на 70 строк, где сформулировал четыре ключевых принципа сотрудничества с ИИ. Файл собрал на GitHub больше 70 000 звёзд и стал самым цитируемым сводом правил среди практиков ИИ-программирования по всему миру. Эти четыре принципа — не пустые лозунги, а практические привычки мышления, выведенные из большой практики. На их основе книга добавляет пятый принцип — итеративную проверку — и получает полную систему из пяти принципов.

\begin{tcbinfo}{Почему файл на 70 строк так ценен?}
70 000 звёзд — это много? Подавляющее большинство проектов на GitHub за всю их жизнь не добираются до этой цифры. То, что текстовый файл на 70 строк её получил, доказывает ровно одну вещь: в эпоху ИИ дефицитен не код, а способность «думать ясно». Эти 70 строк — концентрат умения «думать ясно».

\end{tcbinfo}

\booksection{Принцип 1: Think Before Coding (сначала думай, потом кодируй)}

Это самый контринтуитивный из пяти принципов. Раз ИИ пишет код за нас, зачем «сначала думать»? Ответ прост: качество результата ИИ напрямую зависит от качества ваших входных данных. Если вы сами не разобрались, что хотите сделать, ответы ИИ лишь затянут вас в бесконечный цикл «правим и правим» — в народе это зовётся «ИИ заблудился вместе с вами».

«Сначала думать» значит: прежде чем давать ИИ инструкцию, ответьте себе на три вопроса, образующие неразрывный «мысленный треугольник»:

\begin{enumerate}
  \item Какую проблему я решаю? (цель)
  \item Как выглядит желаемый результат? (критерии приёмки)
  \item Какие ограничения нужно соблюсти? (границы)
\end{enumerate}

\begin{center}
\begin{tabularx}{\linewidth}{|X|X|X|}
\hline
\textbf{Аспект мышления} & \textbf{Размытый ввод (плохо)} & \textbf{Ясный ввод (хорошо)} \tabularnewline
\hline
Определение задачи & Сделай мне сайт & Сделай мне сайт-инструмент для личного пользования с карточками-входами к трём небольшим ИИ-утилитам \tabularnewline
Ожидаемый результат & Чтобы выглядело покрасивее & Адаптивная вёрстка, тёмная тема, анимация карточек при наведении, работает на мобильных \tabularnewline
Ограничения & (не указано) & Без фреймворков, чистые HTML+CSS+JS, деплой одним файлом, бэкенд не нужен \tabularnewline
\hline
\end{tabularx}
\end{center}

\begin{center}\small\color{gray}Сравнение размытого и ясного ввода для одной и той же задачи\end{center}

\begin{tcbwarn}{Внимание: сначала думать $\neq$ сначала писать идеальный документ}
«Сначала думать» не значит «сперва написать идеальный PRD». Это быстрый, но упорядоченный процесс размышления: 30 секунд на три вопроса выше эффективнее, чем 30 минут правок в диалоге. Для мелкой задачи три вопроса могут быть совсем короткими; для крупного проекта их стоит развернуть в полноценный документ требований (подробнее в главе 40).

\end{tcbwarn}

\booksection{Принцип 2: Simplicity First (простота превыше всего)}

Главная ловушка ИИ-программирования — не «не получается сделать», а «сделано слишком много». Попросите у ИИ «систему управления пользователями» — и он, возможно, сгенерирует целый комплект: аутентификацию, права доступа, роли, журналы аудита, экспорт данных — тогда как вам сейчас может понадобиться всего лишь простая страница входа. Избыточная генерация — основная причина «раздувания» ИИ-кода.

«Простота превыше всего» требует в каждой инструкции следовать принципу «минимально достаточного»: минимальным кодом — ключевую функцию, простейшим решением — задачу, наименьшим набором функций — проверку идеи. У Karpathy есть меткое сравнение: «Чтобы посчитать цену со скидкой, вам не нужен движок электронной коммерции».

На практике «простота превыше всего» означает несколько конкретных рабочих привычек:

\begin{itemize}
  \item Каждый раз просите ИИ реализовать одну функциональную точку, а не описывайте всю систему одним махом;
  \item Предпочитайте стандартные библиотеки и простые решения, а не сложные зависимости;
  \item Если задача решается 50 строками кода, не позволяйте ИИ написать 500;
  \item Когда результат ИИ выходит за рамки ожидаемого, смело говорите ему «отрежь лишнее».
\end{itemize}

\begin{center}
\begin{tabularx}{\linewidth}{|X|X|X|}
\hline
\textbf{Сценарий} & \textbf{Избыточное решение (избегать)} & \textbf{Простое решение (рекомендуется)} \tabularnewline
\hline
Нужно хранить данные & Поднять базу данных MySQL + ORM & Сначала файл JSON или SQLite \tabularnewline
Нужен вход пользователя & OAuth2.0 + JWT + роли и права & Сначала простая проверка пароля \tabularnewline
Нужен фронтенд-фреймворк & React + Redux + Router & Сначала попробуйте чистый HTML + ванильный JS \tabularnewline
Нужен деплой в продакшен & Docker + CI/CD + K8s & Сначала заставьте работать локально \tabularnewline
\hline
\end{tabularx}
\end{center}

\begin{center}\small\color{gray}Сравнение избыточных и простых решений в типовых сценариях\end{center}

\booksection{Принцип 3: Surgical Changes (хирургически точные правки)}

Дорабатывая существующий проект, каждое изменение должно быть точным, как хирургическая операция — трогать только то, что нужно, и не задевать лишнего. В ИИ-программировании это особенно важно, потому что у ИИ есть склонность «заодно подправить»: попросите починить один баг — он по пути отрефакторит полфайла.

Ключевое требование «хирургической точности»: в каждой инструкции для ИИ явно задавайте объём и границы изменений. Говорите: «измени только функцию calculate\_discount», «не трогай остальное», «сохрани текущий стиль кода». Если правки ИИ вышли за запрошенные рамки — немедленно требуйте откатить лишнее.

Karpathy иллюстрировал это на примере diff'а (сравнения изменений). Допустим, вы просите ИИ добавить в функцию расчёта скидки корзины логику «клубной цены»: хороший результат — изменена только эта функция и добавлено несколько строк; плохой — переписан весь файл заново. Пусть новый код и «изящнее», но вы не можете гарантировать, что в нём не появились новые баги. В ИИ-программировании «выглядит лучше» никогда не может стать оправданием отступления от инструкции.

Плохая инструкция (границы размыты)

\begin{lstlisting}
Помоги оптимизировать логику расчёта скидок в этом проекте
\end{lstlisting}

Хорошая инструкция (объём задан точно)

\begin{lstlisting}
Измени только функцию calculate\_discount в cart.py: 1. Поверх имеющихся скидок добавь логику клубной цены 2. Участники клуба получают дополнительную скидку 10\% 3. Не меняй сигнатуру функции и другие функции 4. Сохрани текущий стиль имён переменных
\end{lstlisting}

\booksection{Принцип 4: Goal-Driven Execution (исполнение от цели)}

Этот принцип отвечает на ключевой вопрос: давать ли ИИ инструкцию «как делать» или «что делать»? Ответ Karpathy — второе. Вместо «сначала создай класс, потом добавь метод, потом напиши тесты» лучше сказать: «мне нужен инструмент, который разбирает CSV-файл и считает количество уникальных значений в каждом столбце». Пусть ИИ сам выбирает путь реализации, а вы принимаете результат.

Выгода «исполнения от цели» двойная. ИИ получает больше свободы и может выбрать оптимальный путь реализации, вместо того чтобы быть скованным вашими промежуточными шагами; а вы переносите внимание с «как делать» на «правильно ли сделано» — а это и есть дело командира.

\begin{tcbwarn}{Внимание: исполнение от цели $\neq$ пусть идёт как идёт}
Оно требует от вас строгой приёмки после того, как ИИ отработал: удовлетворяет ли результат требованиям? Обработаны ли граничные случаи? Нет ли очевидного переусложнения? Если результат не соответствует ожиданиям, дайте конкретную обратную связь — не «переделай», а «проблема вот в чём: на пустом списке функция падает с IndexError; добавь, пожалуйста, проверку на пустое значение».

\end{tcbwarn}

\begin{center}
\begin{tabularx}{\linewidth}{|X|X|X|}
\hline
\textbf{Стиль инструкции} & \textbf{Пример} & \textbf{Особенности результата} \tabularnewline
\hline
Императивный (как делать) & «Сначала создай Python-класс с именем DataParser, добавь метод parse, принимающий параметр filename, читай данные через модуль csv…» & ИИ превращается в машинистку, а ваши промежуточные шаги могут оказаться не самым оптимальным путём \tabularnewline
Целевой (что делать) & «Мне нужен инструмент: читает CSV-файл, считает количество уникальных значений в каждом столбце, выводит в формате JSON. Дай полную реализацию.» & ИИ сам выбирает оптимальный путь, вы отвечаете за приёмку результата \tabularnewline
\hline
\end{tabularx}
\end{center}

\begin{center}\small\color{gray}Сравнение императивной и целевой инструкций\end{center}

\booksection{Принцип 5: Iterate and Verify (итерируй и проверяй)}

Это пятый принцип, добавленный книгой к четырём принципам Karpathy. Его базовое убеждение: первая выдача ИИ почти никогда не бывает финальной версией. Эффективное ИИ-программирование — это итеративный цикл «генерация — проверка — исправление», а не разовое «инструкция — результат».

«Итерация и проверка» — не пассивное «не так, пробуем ещё раз», а структурированный процесс проверки. После каждой выдачи ИИ проверяйте результат на трёх уровнях:

\begin{itemize}
  \item Уровень 1. Проверка функциональности: оно вообще запускается? Ввод и вывод корректны?
  \item Уровень 2. Проверка границ: что будет с пустым вводом, слишком длинным вводом, некорректным вводом?
  \item Уровень 3. Проверка архитектуры: разумна ли структура кода? Нет ли жёстко вшитых значений? Удобно ли расширять дальше?
\end{itemize}

Определить цель

ИИ пишет код

Проверка функционала

Пройдено

Проверка граничных случаев

Пройдено

Проверка архитектуры

Пройдено — переходим к следующему шагу

Не пройдено

Дать конкретную обратную связь

Сгенерировать заново

\begin{center}\small\color{gray}Рисунок 2: цикл итеративной проверки в ИИ-программировании\end{center}

Ключ этого цикла — шаг «дать конкретную обратную связь». Когда проверка не пройдена, не говорите «неверно, переделывай», а скажите: «функция падает с IndexError на пустом списке; добавь проверку на пустоту в строке 15». Конкретная обратная связь позволяет ИИ исправить точно; расплывчатая лишь заставит его заново угадывать ваши намерения.

\booksection{Harness Engineering: Agent = Model + Harness}

К середине 2026 года в сообществе ИИ-программирования завоёвывало популярность понятие Harness Engineering. Его ключевая формула проста:

Возможности агента = возможности модели × качество обвязки

Возможности модели — интеллектуальный потолок «мозга» ИИ; Harness — «рабочая среда», выстроенная вокруг модели: управление контекстом, интеграция инструментов, оркестрация процессов, системы памяти и прочее.

Это понятие объясняет частое недоумение: почему с тем же GPT-4 или Claude одни делают восхитительные проекты, а другие получают посредственный код? Разница не в модели, а в Harness. Хороший Harness включает: ясный файл со спецификацией проекта, разумную интеграцию инструментов, упорядоченный контекст диалога, отлаженный процесс проверки. Модель задаёт потолок интеллекта, а Harness определяет, насколько близко к этому потолку вы сможете подобраться.

\begin{tcbinfo}{Метафора гениального стажёра}
Представьте модель как гениального стажёра: IQ зашкаливает, но если не дать ему документов о проекте, не рассказать, какими инструментами пользоваться, не назначить сроки и критерии приёмки — хорошего он не сделает. Harness Engineering — это инженерия «как обустроить рабочую среду гениального стажёра».

\end{tcbinfo}

Последующие главы книги по сути учат строить хороший Harness: глава 39 — управление контекстом (система памяти проекта), глава 40 — определение требований (PRD), глава 43 — техника диалога (слои промптов), глава 45 — обеспечение качества (процесс проверки). Всё это вместе и есть ваш Harness.

\begin{tcbpractice}{Практика: разложите свою следующую идею по «мысленному треугольнику»}
Выберите маленькую штуку, которую давно хотели сделать, но руки не доходили (например, веб-страницу с обратным отсчётом дней или скрипт-напоминалку попить воды). За 3 минуты ответьте: (1) Какую проблему она решает? (2) Как она будет выглядеть? (3) Какие ограничения (время, технологии, окружение)? Затем запишите ответы одним абзацем. Поздравляем — вы только что впервые в жизни проделали «Think Before Coding»: этот абзац и есть качественная инструкция.

\end{tcbpractice}

\begin{tcbsummary}{Краткие итоги главы}
\begin{itemize}
  \item Пять принципов восходят к 70-строчному CLAUDE.md Karpathy (70 000+ звёзд); пятый — «итерируй и проверяй» — добавлен книгой
  \item Think Before Coding: прежде чем заговорить, ответьте на три вопроса — что делаем? как выглядит результат? какие ограничения?
  \item Simplicity First: главная ловушка ИИ — «сделать слишком много»; для расчёта скидки не нужен движок интернет-магазина
  \item Surgical Changes: правьте точно, как хирург, и отказывайте ИИ в привычке «заодно подправить»
  \item Goal-Driven Execution: говорите «что делать», а не «как делать», — пусть ИИ сам выбирает путь
  \item Iterate and Verify: проверка на трёх уровнях (функциональность/границы/архитектура) + конкретная обратная связь, итерации по кругу
  \item Harness Engineering: результат = модель × Harness — относитесь к ИИ как к «гениальному стажёру» и управляйте его средой
\end{itemize}

\end{tcbsummary}
